%قبل از شروع مسئله ترجیحاً چیزی ننویسید. اما اگر خواستید بنویسید، در قالب لیستی (فقرات و شمارش) بنویسید.
%برای اضافه کردن نمره، بعدی از شروع‌مسئله، [x] را اضافه کنید و به جای x نمره مورد نظر را قرار دهید. برای زیرمسئله‌ها هم به همین صورت است.
%\شروع‌مسئله[10]
\شروع‌مسئله [15]

%\متن‌سیاه{سوال عملی اول:}\\
%\متن‌سیاه{دراپ‌تست}\\
روزی شخصی به صرافی می‌رود و با دیدن نرخ تبدیل ارزها ایده‌ای به ذهنش می‌رسد. او می‌بیند که می‌تواند با هر پوند
$1.2$
دلار، با هر دلار
$0.9$
یورو و با هر یورو
$0.95$
پوند بخرد. به این ترتیب او می‌تواند با تبدیل ۱ پوند به دلار، دلار به یورو و یورو به پوند صاحب
$1.2 \times 0.9 \times 0.95 = 1.026$
پوند شود و 
$0.026$
پوند سود کند.

فرض کنید
$n$
واحد پول و یک ماتریس
$T$
با ابعاد
$n \times n$
داریم که 
$T_{ij}$
یعنی با یک واحد از $c_i$ می‌توان
$T_{ij}$
واحد از
$c_j$
خرید.
\شروع{زیرمسائل}
\زیرمسئله[7]
یک الگوریتم  از	
$O(n^3)$
ارائه کنید که با ورودی گرفتن ماتریس
$T$
مشخص کند آیا دنباله‌ای از واحدهای پول مانند
$c_{i_1}, c_{i_2}, ..., c_{i_k}$
وجود دارد که تبدیل زنجیره‌ای از
$c_{i_1}$ به $c_{i_2}$
و ...
نهایتا از
$c_{i_k}$ به $c_{i_1}$
باعث سود شود.

\زیرمسئله[8]
یک الگوریتم از
$O(n^3)$
ارائه کنید که در صورتی که دنباله‌ای با شرایط بالا وجود داشته باشد، یک دنباله دلخواه با آن شرایط را خروجی دهد.
\پایان{زیرمسائل}
%\متن‌سیاه{ورودی:}\\
%در خط اول ورودی، ابتدا k و سپس m داده می‌شود.\\
%\متن‌سیاه{خروجی:}\\
%در تنها سطر خروجی، تعداد آیفون‌های مورد نیاز می‌آید. اگر در هیچ صورتی چنین کاری ممکن نبود ۱- چاپ می‌شود.\\
%\متن‌سیاه{محدودیت‌ها:}\\
%$1\leq k\leq 10^4$

%$m\leq10^9$
%\\
%\textbf{\متن‌سیاه{ورودی و خروجی نمونه}\\
%% تست اول
%\begin{latin}
% \begin{tabularx}{\textwidth}{|X | X|}
% \hline
% \textbf{stdin} &
% \textbf{stdout}\\
% \hline
% 100 100&
% 1\\
% \hline
% 
% \end{tabularx}
%\end{latin}
%
%% تست دوم
%
%\begin{latin}
% \begin{tabularx}{\textwidth}{|X | X|}
% \hline
% \textbf{stdin} &
% \textbf{stdout}\\
% \hline
% 340 20&
% 3\\
% \hline
% 
% \end{tabularx}
%\end{latin}}
% صورت اصلی مسئله را در این‌جا بنویسید.



\solution{
	آ) برای اینکه یک دور سود آور باشد باید داشته باشیم
	$$T_{i_1i_2}T_{i_2i_3}...T_{i_ki_1} > 1$$
	$$-\log{(T_{i_1i_2})} - \log{(T_{i_2i_3})} - ... - \log{(T_{i_k}T_{i_1})} < 0$$
	گراف جهت‌دار 
	$G$
	را در نظر بگیرید که هر کدام از راس‌ها معادل با یکی از ارزهاست و بین هر دو راس
	$C_i$ و $C_j$ که
	$i \neq j$
	یک یال از
	$C_i$ به $C_j$
	 با وزن
	$-\log{(T_ij)}$
	وحود دارد. حال با استفاده از الگوریتم بلمن فورد با راس شروع دلخواه تشخیص می‌دهیم آیا در این گراف دور با مسیر منفی وجود دارد یا خیر که همان جواب مسئله ماست. با توجه به اینکه این گراف 
	$n$
	راس و
	$n(n-1)$
	یال دارد زمان اجرای الگوریتم بلمن فورد از مرتبه
	$O(n^3)$
	است.
	
	ب) ابتدا گراف
	$G$
	را مانند بخش الف می‌سازیم. سپس یک راس شروع دلخواه انتخاب کرده و مطابق الگوریتم بلمن فورد
	$n-1$
	بار عملیات ریلکس کردن یال‌ها را انجام داده و سپس یک بار دیگر نیز این عملیات را انجام می‌دهیم. اگر در دور آخر فاصله راسی تغییر کرد این راس در یک دور منفی است. حال کافیست هر سری به راس قبلی راس فعلی (که بلمن فورد مشخص کرده) برویم و این کار را آنقدر تکرار کنیم تا یک راس تکراری ببینیم. دور پیدا شده، یک دور منفی است.(چرا؟) هزینه اصلی این الگوریتم $n$ بار ریلکس کردن یال‌هاست بنابراین زمان اجرای این الگوریتم از مرتبه
	$O(n^3)$
	است.
	

	
}

%%%%%%%%%%%%%%%%%%%%%%%%%%%%%%%%%%%%%%%%%%%%%%%%%%%%%
% برخی دستورات مفید برای تمرین‌ها. %
%%%%%%%%%%%%%%%%%%%%%%%%%%%%%%%%%%%%%%%%%%%%%%%%%%%%%

% فرمول‌های نمونه:
%$f(x) = \Theta(x)$
%$T(x, y) \leq \log(x) + 2^y + \binom{2n}{n}$
% {\tt code\_function}


% لیست موارد بدون شماره‌گذاری:
%\شروع{فقرات}
% \فقره مورد اول 
% \شروع{فقرات}
% \فقره مورد اول داخلی 
% \شروع{فقرات}
% \فقره اولین مورد خیلی داخلی 
% \فقره دومین مورد خیلی داخلی 
% \پایان{فقرات}
% \فقره مورد دوم داخلی 
% \پایان{فقرات}
% \فقره مورد دوم 
%\پایان{فقرات}

% لیست موارد شماره‌گذاری شده:
%\شروع{شمارش}
% \فقره یک مورد 
% \فقره یک مورد دیگر 
% \فقره باز هم یک مورد دیگر
%\پایان{شمارش}

% برای نوشتن یک فرمول چند خطی که یک شماره بخوره. برای حالت‌های دیگر به سایت زیر مراجعه کنید:
% https://www.overleaf.com/learn/latex/Aligning_equations_with_amsmath
%\begin{align} 
% \begin{split}
% (x+y)^3 &= (x+y)^2(x+y) \\
% &= (x^2+2xy+y^2)(x+y) \\
% &= (x^3+2x^2y+xy^2) + (x^2y+2xy^2+y^3) \\
% &= x^3+3x^2y+3xy^2+y^3
% \end{split} 
%\end{align}

% نوشتن ماتریس و آرایه دوبعدی
%\begin{align}
% A = 
% \begin{bmatrix}
% A_{11} & A_{21} \\
% A_{21} & A_{22}
% \end{bmatrix}
%\end{align}

% اضافه کردن تصویر با شرح. برای اضافه کردن تصویر بدون شرح، فقط دستور درج‌تصویر را استفاده کنید.
%\شروع{شکل}[h!]
%\تنظیم‌ازوسط
% \درج‌تصویر[پهنا=0.5\پهنای‌متن]{img.jpg}
%\شرح{شرح تصویر}
%\برچسب{تصویر:۱}
%\پایان{شکل}
%\پایان{زیرمسائل}

%کد زیر را ببینید.
%%\lstinputlisting{source_filename.py}
%\begin{latin}
%\begin{lstlisting}
%# Import the modules
%import sys
%import random
%
%ans = True
%
%while ans:
%question = raw_input("Ask the magic 8 ball a question: (press enter to quit) ")
%
%answers = random.randint(1,8)
%
%if question == "":
%sys.exit()
%
%elif answers == 1:
%print "It is certain"
%
%elif answers == 2:
%print "Outlook good"
%
%elif answers == 3:
%print "You may rely on it"
%
%elif answers == 4:
%print "Ask again later"
%
%elif answers == 5:
%print "Concentrate and ask again"
%
%elif answers == 6:
%print "Reply hazy, try again"
%
%elif answers == 7:
%print "My reply is no"
%
%elif answers == 8:
%print "My sources say no"
%\end{lstlisting}
%\end{latin}